% Last Updated: March 28, 2023
	\documentclass[a4paper, twoside]{article}
	% Margins
	\usepackage{geometry}
		\geometry{
			% showframe,
			left   = 20mm,
			right  = 20mm,
			top    = 20mm,
			bottom = 20mm
		}

	% Headers and footers
	\usepackage{fancyhdr}
		\pagestyle{fancy}
		\fancyhf{}
		\fancyhead[LE, RO]{\thepage}
		\fancyhead[CE]{\textsc{Cong Wen}}
		\fancyhead[CO]{\textsc\rightmark}

		\setlength{\headheight}{15pt}
		\renewcommand\headrulewidth{0pt}








\input{../universal-pre}

\title{\tbf{Learning note of AWS}}
\author{\tbf{Notes} by Cong Wen, 07/2023}
\date{}
% Last Updated: March 28, 2023



\begin{document}
\maketitle
\thispagestyle{empty}
\begin{abstract}
	This is my self-use notes\footnote{Last updated on July 24, 2023} when reading materials of \tbf{Arizona Winter School}. Please reach out if you find anything incorrect.

\end{abstract}
\tableofcontents
% ----------------------------------------------------------------------------------------------------


\sct{2021-Chan-Quadratic forms}


\ssc{Results}

The content covered classification theorems of quadratic forms over $K$ with $\Char{K}\neq2$ and applications.


Case $K=\bF_{q}$, they are determined: dimension, discriminant.

Case $K=\bQ_{p}$, they are determined: dimension, discriminant, Hasse invariant.

Case $K=\bR$, they are determined: signature.

Case $K=\bQ$, they are determined: restrictions to $\bQ_{p}$ for $p\lqs\ift$ (Hasse principle holds).


\ssc{Applications}

\begin{thm}
	Let $n\in\bZ_{>0}$, then
	\begin{itm}
		\item (Gau{\ss}) $n$ is a sum of three squares if and only if $n\neq4^{a}(8b-1),\fal a,b\in\bZ$.
		\item (Gau{\ss}) $n$ is a sum of three triangular numbers.
		\item (Lagrange) $n$ is a sum of four squares.
	\end{itm}
\end{thm}


$\Box, \km_{\Sgm}^{\Box,2}$ is good.

nice nice

% ----------------------------------------------------------------------------------------------------
\bibliographystyle{amsalpha}
\bibliography{../universal-ref}
\end{document}

